% Preface: the README's About, Verification and Exercises sections, and the course info's
% prerequisites, as the front matter of a book.

\chapter{Preface}
\markboth{Preface}{Preface}

This book is for embedded engineers who write C or C++ for microcontrollers and have never
written assembly, or have read some and never produced any. It takes you from your first
\code{ldi} to a driver library with interrupts, timers and a watchdog, written entirely in AVR
assembly and verified by tests that run on your laptop.

\textbf{The subject is the machine.} Assembly is how you get at it; it is not being taught here as
a way to write applications, and this book will not tell you it is faster than your compiler.
What it is good for is narrower and more durable:
\begin{itemize}
  \item Reading a disassembly and knowing what you are looking at.
  \item Writing the handful of things C genuinely cannot express.
  \item Understanding what your compiler was doing all along.
\end{itemize}

\textbf{One body of code is built across the book, and it is yours to write:} an AVR assembly
driver library, grown one chapter at a time. \code{led.asm}, \code{btn.asm},
\code{led_array.asm}, \code{timer.asm}, \code{watchdog.asm} and the utilities under them, plus an
application program that composes them.

\section*{What you are assumed to know}
The subject of this book is \textbf{AVR assembly and the machine underneath it}, not embedded
programming in general. You are expected to arrive already comfortable with:
\begin{itemize}
  \item C or C++ for microcontrollers: pointers, structs, bitwise operators, \code{volatile}.
  \item Binary and hexadecimal, two's complement, and reading a bit mask such as \code{(1 << 5)}.
  \item What a microcontroller peripheral is, and the idea that a register write configures
    hardware.
  \item Reading a datasheet section: finding a register, its bit fields, and their meanings.
\end{itemize}
None of that is taught from scratch. What the book does teach from nothing is the instruction set,
the register file, the calling contract, the vector table, and the toolchain: \code{avra} as an
assembler, its listing and map files, and simavr as the thing that tells you whether you were
right.

\textbf{No assembly experience of any kind is assumed}, on the AVR or anywhere else. If you have
written x86 or ARM assembly, Chapter~1 will be short work and the rest will not be: the AVR's
16-bit-pointer-in-three-register-pairs arrangement, its separate program and data address spaces,
and its I/O register window are all different enough to be worth reading rather than skimming.

\section*{How verification works}
Every quantitative claim in this book is checked twice, and the whole method is built on the two
not always agreeing:
\begin{itemize}
  \item \textbf{By hand.} You compute the number from the instruction set summary and the
    datasheet, before any code exists. Nothing automated checks this step, deliberately, and that
    is the point of it.
  \item \textbf{In the simulator.} The assembly you wrote is assembled, loaded into simavr by a
    harness the companion repository provides, and measured: cycles consumed, pins driven, SRAM
    written, interrupts taken.
\end{itemize}
Where the two disagree, the reconciliation is the lesson, and every chapter has one exercise that
is exactly this. Your hand count of \code{shift_bits} and the figure the simulator reports at a
call site differ by a fixed amount, and the difference is the \code{ldi} that set the argument up
and the \code{rcall} that got you there. Your timer's computed period at 1\,kHz and its
measured period agree, and then at 3\,kHz they do not, and the gap is not the timer's error but
the two-cycle instruction each interrupt has to wait for.

And once, in Chapter~3, they disagree because the \textbf{simulator} is wrong: simavr takes an
interrupt without charging the four cycles the datasheet says the hardware spends pushing the
program counter. Knowing which of your two instruments to believe, and why, is the skill this
arrangement exists to build.

\textbf{No exercise in this book needs hardware.} Everything is written, assembled, simulated and
verified on an ordinary laptop. If you own an Arduino Uno or Nano with an ATmega328P and want to
watch a real LED, \secref{c2:sec:board} covers flashing it with \code{avrdude}; nothing depends
on it.

\section*{What the exercises look like}
Every chapter ends with six to ten exercises, and each one is labelled with its kind:

\begin{booktable}{@{}p{9.4em}L@{}}
  \thead{Kind} & \thead{What it asks of you} \\
  \midrule
  \kindtag{Recall} & State something from the chapter, in your own words. No tools. \\
  \kindtag{Hand calculation} & Work a number out on paper, from the datasheet and the instruction
    set summary. \\
  \kindtag{Design} & Decide something before writing it: a register allocation, a struct layout,
    a prescaler. \\
  \kindtag{Code} & Write assembly. Checked by that chapter's test suite, in the simulator. \\
  \kindtag[fillaccent]{Cross-check} & Exactly one per chapter. See below. \\
\end{booktable}

\textbf{The Cross-check is the signature exercise of this book.} You compute a number by hand,
you run your own code on the same problem, and you reconcile the two. Its solution says how large
a discrepancy to expect and where it comes from, because ``your two answers differ by 8\%'' is the
lesson, not a sign that you did it wrong. Some of these discrepancies close to zero and some do
not, and knowing which is which is most of what separates an engineer who can time a routine from
one who assumes the datasheet did it for them.

\textbf{The worked solutions are not in this book.} They will be published in the companion
repository, beside the lecture each chapter is typeset from, and until then most of the code exercises need nothing more than
the chapter's test suite to tell you whether you are right. The hand calculations have short
answers in Appendix~\ref{app:answers}, numbers only, so that you can check one without a laptop.
Clone the repository before you start; the book assumes you have it open.

\section*{Setting up}
Everything runs on WSL or Ubuntu. Clone the companion repository with its two submodules, the
test framework and the simulator harness, at the tag this edition was written against, and install
the toolchain. The tag is \code{v1.0.0}: the code and the test suites exactly as this book
describes them, whatever the main branch has done since.

\begin{shell}
git clone --recursive --branch v1.0.0 \
          https://github.com/qrtech-academy/embedded-asm.git
sudo apt -y update
sudo apt -y install git make g++ avra binutils-avr simavr \
                    libsimavr-dev libelf-dev clang-format
sudo apt -y install gcc-avr avr-libc  # Chapter 6 only
avra --version                        # 1.4.2 or newer
\end{shell}

\noindent Four commands then cover everything the book asks of the repository:

\begin{shell}
make help                # List every target.
make build               # Assemble the drivers and build the simulator harness.
make test                # Build and run every chapter's test suite.
make clean               # Remove everything the build generates.
\end{shell}

\noindent Every one of those runs from the very first commit, when almost nothing exists, and
reports what it could not do by name. A check that silently did not run must never read like one
that passed, so \code{make build} on a fresh clone builds the two things that ship, prints a
\code{SKIP} line for the drivers you have not written, and exits 0. That is the truth about a
repository where you have not written anything yet.
